\chapter{Mean-Field Viewpoint}
\label{cha:mean-field-appendix}

The main text uses the NTK viewpoint as its main analytical framework because
it gives a fixed operator on functions. On \(S^1\), this operator can be
diagonalized in the Fourier basis, which makes it possible to discuss
frequency-dependent learning in terms of eigenvalues and eigenspaces.

The mean-field viewpoint gives another way to take a large-width limit. It is
not used in the main finite-sample or finite-width analysis, but it is useful
context because it describes a different aspect of wide-network training. In
the mean-field scaling, one does not primarily track a fixed tangent kernel at
initialization. Instead, one tracks how the hidden neurons are distributed in
parameter space, and how this distribution changes during training.

This appendix records the corresponding derivation for a two-layer network. The
purpose is to keep the comparison available without interrupting the main line
of the thesis.

\section{\label{app:mf-overview}Overview and relation to the NTK viewpoint}

Consider a two-layer network with the mean-field scaling
\begin{equation}
	f_\theta(x)
	=
	\frac{1}{m}
	\sum_{\alpha=1}^m
	v_\alpha\,\varphi(w_\alpha^\top x).
	\label{eq:app-mf-network-overview}
\end{equation}
Here each neuron has parameters
\[
	\theta_\alpha=(v_\alpha,w_\alpha)
	\in
	\Omega:=\mathbb R\times\mathbb R^d .
\]
The key observation is that the network output depends on the collection of
neurons only through their empirical distribution,
\[
	\mu_m
	:=
	\frac{1}{m}
	\sum_{\alpha=1}^m \delta_{\theta_\alpha}.
\]
Indeed,
\[
	f_\theta(x)
	=
	\int_\Omega v\,\varphi(w^\top x)\,d\mu_m(v,w).
\]
Thus the state of the network can be described by a probability measure over
neuron parameters.

In suitable large-width limits, the empirical measure \(\mu_m\) converges to a
deterministic time-dependent measure \(\mu_t\). Training is then described by
the evolution of this measure. This viewpoint was developed for two-layer
networks by \citet{Mei2018meanfieldview}, and related interacting-particle
descriptions were studied by \citet{Rotskoff_2022}. From an optimization
perspective, \citet{chizat2018globalconvergencegradientdescent} studied the
corresponding many-particle limit using tools from optimal transport.
Extensions and rigorous frameworks for deeper networks have also been studied
\citep{nguyen2019meanfieldlimitlearning,araujo2019meanfieldlimitcertaindeep,
	sirignano2019meanfieldanalysisneural,nguyen2023rigorousframeworkmeanfield}.

This viewpoint is useful because it can describe feature movement in the
large-width limit. In contrast, the strict NTK limit freezes the tangent kernel:
the features defined by the gradients at initialization remain fixed in the
limit. For the spectral question studied in this thesis, however, the NTK
viewpoint is more directly useful. The main object here is an operator on
functions whose eigenvalues and eigenspaces can be compared with Fourier modes
on \(S^1\). The mean-field viewpoint instead describes the evolution of a
probability measure over parameters, so it does not give a single fixed
Fourier-diagonal operator in the same way.

\section{\label{app:mf-detailed}Detailed derivation of the mean-field representation}

\section{\label{app:mf-detailed}Detailed derivation of the mean-field representation}

\subsection{\label{app:mf-empirical-measure}Parameter space and empirical measure}

Let $\Omega := \mathbb{R} \times \mathbb{R}^d$
denote the parameter space, and let $\mathcal{F}$ be its Borel $\sigma$-algebra.
For each neuron $\alpha$, define the Dirac measure $\delta_{\theta_\alpha}$ on
$(\Omega,\mathcal{F})$, i.e. if $A \in \mathcal{F}$ then,
\[
	\delta_{\theta_\alpha}(A)
	=
	\begin{cases}
		1, & \theta_\alpha \in A, \\
		0, & \text{otherwise}.
	\end{cases}
\]
Define the empirical neuronal measure
\begin{equation}
	\mu_m := \frac{1}{m}\sum_{\alpha=1}^m \delta_{\theta_\alpha}.
	\label{eq:app-empirical-measure}
\end{equation}
By construction, $\mu_m \ge 0$ and  $\mu_m(\Omega)=1$, hence $\mu_m \in
	\mathcal P(\Omega)$ where $\mathcal P(\Omega)$ is the space of Borel
probability measures on $\Omega$.

\paragraph{Integration against Dirac measures.}
Let $\chi_A$ denote the indicator function of a measurable set $A \subset \Omega$.
Then, for any $\theta_\alpha \in \Omega$,
\[
	\int_\Omega \chi_A(\theta)\,d\delta_{\theta_\alpha}(\theta)
	=
	\delta_{\theta_\alpha}(A)
	=
	\chi_A(\theta_\alpha).
\]
More generally, if $s(\theta) = \sum_{k=1}^r c_k \chi_{A_k}(\theta)$ is a simple
function, then
\[
	\int_\Omega s(\theta)\,d\delta_{\theta_\alpha}(\theta)
	=
	\sum_{k=1}^r c_k \chi_{A_k}(\theta_\alpha)
	=
	s(\theta_\alpha).
\]
Any nonnegative measurable function $g : \Omega \to \mathbb{R}$ can be approximated by
a sequence of simple functions, $\{s_n\}_{n\ge 0}$. If $s_n \uparrow g$ and each $s_n \ge 0$ then
by the Monotone Convergence Theorem,
\begin{equation}
	\int_\Omega g\,d\delta_{\theta_\alpha}
	= \int_\Omega \lim_{n \to \infty} s_n\,d\delta_{\theta_\alpha}
	=  \lim_{n \to \infty} \int_\Omega s_n\,d\delta_{\theta_\alpha}
	=  \lim_{n \to \infty} s_n(\theta_\alpha)
	= g(\theta_\alpha).
	\label{eq:dirac-integration}
\end{equation}

\paragraph{Integral representation of the network.}
Consider now the integral of $g$ with respect to the empirical measure $\mu_m$:
\[
	\int_\Omega g\,d\mu_m
	=
	\int_\Omega g\,d\left(\frac{1}{m}\sum_{\alpha=1}^m \delta_{\theta_\alpha}\right)
	=
	\frac{1}{m}\sum_{\alpha=1}^m g(\theta_\alpha).
\]
Define
\[
	g(\theta_\alpha) = g(v_\alpha,w_\alpha) := v_\alpha\,\varphi(w_\alpha^\top x).
\]
Then
\begin{equation}
	\int_\Omega g\,d\mu_m
	=
	\frac{1}{m}\sum_{\alpha=1}^m v_\alpha\,\varphi(w_\alpha^\top x)
	=
	f_\theta(x),
	\label{eq:app-mf-integral-form}
\end{equation}
recovering the finite-width network.

\subsection{\label{app:mf-weak-convergence}Initialization and weak convergence}

Assume that the parameters $\{\theta_\alpha(0)\}_{\alpha=1}^m$ are initialized
i.i.d.\ according to a probability measure $\mu_0$ on $\Omega$.
Define the empirical measure at initialization
\[
	\mu_{m,0}
	:=
	\frac1m\sum_{\alpha=1}^m \delta_{\theta_\alpha(0)} .
\]
Then $\mu_{m,0}$ converges weakly to $\mu_0$ almost surely as $m\to\infty$.
Indeed, for any bounded continuous test function $\psi:\Omega\to\mathbb R$,
\[
	\int_\Omega \psi\,d\mu_{m,0}
	=
	\frac1m\sum_{\alpha=1}^m \psi(\theta_\alpha(0))
	\;\xrightarrow[m \to \infty]{\text{a.s.}}\;
	\mathbb E_{\theta\sim\mu_0}[\psi(\theta)]
	=
	\int_\Omega \psi\,d\mu_0,
\]
where the convergence follows from the strong law of large numbers applied to
the i.i.d.\ random variables $\psi(\theta_\alpha(0))$.
Since this holds for all bounded continuous $\psi$, we conclude that
$\mu_{m,0} \to \mu_0$ \citep{vadaranjan1958convergenceofprobabilitydistribution}.
As a consequence, the network output at initialization converges pointwise to
the deterministic limit
\[
	f_{\mu_0}(x)
	:=
	\int_\Omega v\,\varphi(w^\top x)\,d\mu_0(v,w).
\]
For instance, under i.i.d.\ Gaussian initialization of the parameters, i.e.
$v_0 \sim \mathcal{N}(0,I_m)\; W_0 \sim \mathcal{N}(0, I_m \otimes I_d)$, the
limiting initial measure $\mu_0 = \mathcal{N}(0, I_{d+1})$ is a Gaussian
measure on a $d+1$-dimensional space.

\subsection{\label{app:mf-particle-dynamics}Training dynamics and evolution of the empirical measure}

Let $X=(x_1,\dots,x_n)\in\mathbb R^{d\times n}$ and
$y=(y_1,\ldots,y_n)\in\mathbb R^n$ denote the training set, and consider the
squared loss
\[
	L(v,W)
	=
	\frac12\sum_{i=1}^n\bigl(f_{v,W}(x_i)-y_i\bigr)^2
	=
	\frac12\|f_{v,W}(X)-y\|_2^2,
\]
where $f_{v,W}(X):=(f_{v,W}(x_i))_{i=1}^n$.

We study continuous-time gradient flow for a finite-width network with $m$
neurons and learning rate $\eta>0$,
\[
	\dot v_{t,\alpha}
	=
	-\eta\,\nabla_{v_\alpha}L(v_t,W_t),
	\qquad
	\dot w_{t,\alpha}
	=
	-\eta\,\nabla_{w_\alpha}L(v_t,W_t),
	\qquad
	\alpha\in\{1,\dots,m\}.
\]

A direct computation yields, for each $\alpha$,
\begin{align}
	\dot v_{t, \alpha}
	 & =
	\frac{\eta}{m}\,
	\varphi(X^\top w_{t,\alpha})^\top
	\bigl(y-f_{v_t,W_t}(X)\bigr),
	\label{eq:mf-ode-v-app} \\
	\dot w_{t,\alpha}
	 & =
	\frac{\eta}{m}\,
	X\Bigl(
	\varphi'(X^\top w_{t,\alpha})
	\odot
	\bigl(y-f_{v_t,W_t}(X)\bigr)
	\Bigr)v_{t,\alpha},
	\label{eq:mf-ode-w-app}
\end{align}
where $\odot$ denotes elementwise (Hadamard) multiplication.

As seen in \eqref{eq:mf-ode-v-app}--\eqref{eq:mf-ode-w-app}, evolution of
each neuron depends on the parameters only through the current empirical
measure $\mu_{m,t}$ via $f_{v_t,W_t}(X)$. Thus, the particle system
\eqref{eq:mf-ode-v-app}--\eqref{eq:mf-ode-w-app} induces an evolution of the
empirical measure in parameter space \citep{Mei2018meanfieldview,Rotskoff_2022}.

Hence, the particle system \eqref{eq:mf-ode-v-app}--\eqref{eq:mf-ode-w-app}
can be viewed as an interacting particle system driven by a measure-dependent
velocity field. For any $\theta=(v,w)\in\Omega$ and any $\mu\in\mathcal
	P(\Omega)$, define
\[
	b((v,w);\mu)
	:=
	\begin{pmatrix}
		\varphi(X^\top w)^\top\bigl(y-f_\mu(X)\bigr) \\[2mm]
		X\Bigl(\varphi'(X^\top w)\odot\bigl(y-f_\mu(X)\bigr)\Bigr)\,v
	\end{pmatrix},
\]
where
\[
	f_\mu(x):=\int_\Omega v\,\varphi(w^\top x)\,d\mu(v,w),
	\qquad
	f_\mu(X):=\bigl(f_\mu(x_i)\bigr)_{i=1}^n.
\]
Then the gradient-flow dynamics can be written compactly as
\[
	\dot\theta_{t,\alpha}
	=
	\frac{\eta}{m}\,b(\theta_{t,\alpha};\mu_{m,t}),
	\qquad \alpha=1,\dots,m.
\]
In particular, each particle interacts with the rest of the system only through the
current empirical measure $\mu_{m,t}$ (via $f_{\mu_{m,t}}(X)$).

\subsection{\label{app:mf-continuity-equation}Continuity equation and push-forward formulation}

Equivalently, $\mu_{m,t}$ solves the continuity equation
\begin{equation}
	\partial_t\mu_{m,t}
	+
	\nabla_\theta\!\cdot\!\Bigl(\mu_{m,t}\,\frac{\eta}{m}\,b(\cdot;\mu_{m,t})\Bigr)
	=
	0,
	\label{eq:mf-continuity-compact-app}
\end{equation}
in the sense of distributions, i.e.\ for every test function
$\psi\in C_c^\infty(\Omega)$,
\[
	\frac{d}{dt}\int_\Omega \psi(\theta)\,d\mu_{m,t}(\theta)
	=
	\frac{\eta}{m}\int_\Omega \nabla_\theta\psi(\theta)\cdot b(\theta;\mu_{m,t})\,d\mu_{m,t}(\theta).
\]

Choosing the mean-field scaling $\eta=m$ yields a nontrivial $\mathcal O(1)$
evolution in time and, as $m\to\infty$, one expects $\mu_{m,t} \to \mu_t$,
where the limit $\mu_t$ satisfies
\begin{equation}
	\partial_t\mu_t+\nabla_\theta\!\cdot\bigl(\mu_t\,b(\cdot;\mu_t)\bigr)=0,
	\qquad \mu_0=\mathcal N (0,I_{d+1}).
	\label{eq:mf-continuity-with-initial-conditions-app}
\end{equation}

\paragraph{Push-forward formulation.}
Recall that if $g:\mathcal X\to\mathcal Z$ is measurable and $\mu\in\mathcal P(\mathcal X)$, then the
\emph{push-forward} of $\mu$ by $g$ is the measure $g_*\mu\in\mathcal P(\mathcal Z)$ defined by
\[
	(g_*\mu)(A):=\mu\bigl(g^{-1}(A)\bigr),
	\qquad A\subset\mathcal Z \text{ measurable}.
\]
Let $\Phi_t:\Omega\to\Omega$ denote the flow map associated with the velocity field
$\frac{\eta}{m}\,b(\cdot;\mu_{m,t})$. Then the empirical measure is transported by this flow:
\[
	\mu_{m,t} = (\Phi_t)_*\mu_{m,0}.
\]
In particular, for any measurable observable $h:\Omega\to\mathbb R$,
\[
	\int_\Omega h(\theta)\,d\mu_{m,t}(\theta)
	=
	\int_\Omega h\!\left(\Phi_t(\theta)\right)\,d\mu_{m,0}(\theta),
\]
so evaluating the network at time $t$ is obtained by taking
$h_x(\theta):=v\,\varphi(w^\top x)$ and writing
\[
	f_{\mu_{m,t}}(x)
	=
	\int_\Omega v\,\varphi(w^\top x)\,d\mu_{m,t}(v,w)
	=
	\int_\Omega v\,\varphi(w^\top x)\,d\bigl((\Phi_t)_*\mu_{m,0}\bigr)(v,w).
\]

\section{\label{app:mf-comparison-ntk}Comparison with the NTK scaling}

The distinction between the mean-field and NTK viewpoints can also be seen from
the tangent kernel. Recall that, in the NTK parameterization used in the main
text, the two-layer network is scaled by \(1/\sqrt m\). This scaling gives a
nontrivial limiting tangent kernel at initialization, and in the infinite-width
NTK regime this kernel remains fixed during training.

In the mean-field parameterization,
\[
	f_\theta(x)
	=
	\frac1m
	\sum_{\alpha=1}^m
	v_\alpha\varphi(w_\alpha^\top x),
\]
the empirical tangent kernel is instead
\begin{equation}
	\Theta^{(m)}_t(x,x')
	=
	\frac{1}{m^2}\sum_{\alpha=1}^m
	\Bigl(
	\varphi(w_{t,\alpha}^\top x)\,\varphi(w_{t,\alpha}^\top x')
	+
	v_{t,\alpha}^2
	\varphi'(w_{t,\alpha}^\top x)
	\varphi'(w_{t,\alpha}^\top x')
	x^\top x'
	\Bigr).
	\label{eq:app-mf-ntk-equation}
\end{equation}
This kernel is of order \(m^{-1}\). Under the mean-field time scaling, the
effective kernel driving the output dynamics is the scaled quantity
\(m\Theta^{(m)}_t\). As \(m\to\infty\), this has the formal limit
\[
	\lim_{m\to\infty}m\Theta^{(m)}_t(x,x')
	=
	\int_\Omega
	\Bigl(
	\varphi(w^\top x)\varphi(w^\top x')
	+
	v^2\varphi'(w^\top x)\varphi'(w^\top x')x^\top x'
	\Bigr)
	\,d\mu_t(v,w).
\]
Because the measure \(\mu_t\) changes with time, this limiting kernel also
changes with time. This is the sense in which the mean-field viewpoint can
retain feature movement in the large-width limit.

This is useful for studying representation change, but it is less convenient
for the spectral analysis in the main text. There, the central object is a
fixed operator on \(L^2(S^1)\), whose eigenfunctions and eigenvalues can be
compared with Fourier modes. In the mean-field viewpoint, the limiting object is
the evolving measure \(\mu_t\), and the associated kernel changes with time.
Thus the same fixed Fourier-diagonal picture is not available directly.

There are also technical limitations. The long-time behavior of the limiting
measure-valued dynamics is generally harder to characterize: one may need to
ask whether \(\mu_t\) converges as \(t\to\infty\), and which limiting measure is
selected. Moreover, the simple empirical-measure description above is most
natural for two-layer networks. For deeper architectures, one needs more
careful mean-field constructions, as studied for example by
\citet{nguyen2019meanfieldlimitlearning},
\citet{araujo2019meanfieldlimitcertaindeep},
\citet{sirignano2019meanfieldanalysisneural}, and
\citet{nguyen2023rigorousframeworkmeanfield}.
